%-------------------------------------- COMIENZA descripción del caso de uso.
\begin{UseCase}{CU-09}{Gestión de personal por Coordinador}{
	Permite al Coordinador de Caso registrar nuevos especialistas operativos, editar sus datos personales o laborales, y habilitar o deshabilitar sus cuentas de acceso en el sistema, dentro del ámbito de su jurisdicción sobre el equipo multidisciplinario.
}
	\UCitem{Versión}{\color{Gray}1.0}
	\UCitem{Autor}{\color{Gray}Ramírez Cuevas Emiliano}
	\UCitem{Supervisa}{\color{Gray}Guerra Salinas Edgar Rafael}
	\UCitem{Actor}{\hyperlink{tCoordinador}{Coordinador de Caso}}
	\UCitem{Propósito}{Delegar al Coordinador de Caso la capacidad de administrar las cuentas de los especialistas operativos de la Fundación de forma supervisada.}
	\UCitem{Entradas}{CURP, RFC, Nombre, Apellidos, Fecha de nacimiento, Sexo, Correo, Teléfono, Domicilio, Rol, Cédula profesional, Especialidad y Contraseña provisional del nuevo usuario.}
	\UCitem{Origen}{Teclado}
	\UCitem{Salidas}{Confirmación del registro, edición o cambio de estado del usuario en el panel del Coordinador.}
	\UCitem{Destino}{Pantalla y base de datos relacional.}
	\UCitem{Precondiciones}{El Coordinador debe haber iniciado sesión con credenciales válidas en la plataforma.}
	\UCitem{Postcondiciones}{El personal dado de alta o modificado queda disponible para la formación o ajuste de equipos multidisciplinarios.}
	\UCitem{Errores}{
		\begin{description}
			\item[E1:] Datos obligatorios incompletos o en formato incorrecto (ej. CURP inválida).
			\item[E2:] Correo electrónico o CURP duplicados en el sistema.
		\end{description}
	}
	\UCitem{Tipo}{Caso de uso primario}
	\UCitem{Observaciones}{Regido por las reglas de negocio \hyperlink{BR-06}{BR-06}, \hyperlink{BR-07}{BR-07} y \hyperlink{BR-08}{BR-08}.}
\end{UseCase}

\begin{UCtrayectoria}
	\UCpaso[\UCactor] Accede al panel de administración de personal de Taimotla en el módulo del Coordinador.
	\UCpaso Despliega el \IUref{IU31}{Panel de Control del Coordinador} con la lista de empleados registrados (activos e inactivos), agrupada por especialidad.
	\UCpaso[\UCactor] Realiza una de las siguientes acciones:
		\begin{itemize}
			\item \textbf{Alta:} Presiona el botón \IUbutton{Registrar Nuevo Personal}.
			\item \textbf{Edición:} Presiona el botón \IUbutton{Editar} en el renglón del empleado.
			\item \textbf{Deshabilitar:} Presiona el botón \IUbutton{Deshabilitar} en el renglón del empleado.
			\item \textbf{Habilitar:} Presiona \IUbutton{Habilitar Cuenta} en la sección de Inactivos.
		\end{itemize}
	\UCpaso Despliega el formulario de registro, la pantalla de edición o ejecuta la acción de cambio de estado según corresponda.
	\UCpaso[\UCactor] Llena los campos del formulario y presiona \IUbutton{Guardar Registro Completo}. \Trayref{A}.
	\UCpaso Valida los campos obligatorios del empleado según su rol, incluyendo la cédula profesional para especialistas conforme a la regla \hyperlink{BR-08}{BR-08}. \Trayref{A}.
	\UCpaso Verifica que el correo, RFC y CURP no estén duplicados y cumplan el formato requerido conforme a las reglas \hyperlink{BR-06}{BR-06} y \hyperlink{BR-07}{BR-07}. \Trayref{B}.
	\UCpaso Encripta la contraseña provisional de acceso.
	\UCpaso Guarda la información en la base de datos vinculando al empleado con su rol y especialidad.
	\UCpaso Despliega el mensaje de éxito {\bf MSG-09} confirmando la operación realizada.
	\UCpaso[] -- -- Fin del caso de uso.
\end{UCtrayectoria}

\begin{UCtrayectoriaA}{A}{Datos obligatorios incompletos o con formato inválido}
	\UCpaso Muestra el mensaje de error {\bf MSG-04} indicando los campos vacíos o con formato inválido.
	\UCpaso[\UCactor] Oprime el botón \IUbutton{Aceptar}.
	\UCpaso Regresa al paso 5 de la trayectoria principal.
\end{UCtrayectoriaA}

\begin{UCtrayectoriaA}{B}{Duplicidad de datos de identidad en el sistema}
	\UCpaso Muestra el mensaje de error {\bf MSG-02} alertando que la CURP o el correo electrónico ya existen en otro registro de cuenta activa.
	\UCpaso[\UCactor] Oprime el botón \IUbutton{Aceptar}.
	\UCpaso Regresa al paso 5 de la trayectoria principal.
\end{UCtrayectoriaA}
%-------------------------------------- TERMINA descripción del caso de uso.